% Part: turing-machines 
% Chapter: undecidability
% Section: introduction

\documentclass[../../../include/open-logic-section]{subfiles}

\begin{document}

\olfileid{tur}{und}{int} 
\olsection{Pambuka}

Mbokmenawa katon cetha yen ora saben fungsi, malah
ora saben fungsi aritmetika, bisa diitung. Cacahe
fungsi kakehan lan tumindake ana kang rumit banget.
Contone, fungsi kang ditegesi saka peluruhan partikel
radioaktif utawa saka tumindak liya kang kacau utawa
acak. Upamane kita wiwit ngitung interval-interval
siji detik saka sawijining wektu, banjur netepake
fungsi $f(n)$ minangka cacahe partikel ing jagad raya
kang luruh sajrone interval siji detik kapangka-$n$
sawise wektu wiwitan iku. Fungsi iki katon kaya
tuladha fungsi kang ora bakal bisa kita itung.

Nanging ora bisa mbayangake carane ngitung fungsi
kaya mangkono iku beda karo tenan mbuktekake yen
fungsi iku ora bisa diitung. Malah fungsi kang katon
rumit tanpa pangarep-arep bisa uga, kanthi teges
abstrak, bisa diitung. Contone, upamane umur jagad
raya winates---ing sawijining dina kang isih adoh
banget, jagad raya bakal ngempet dadi siji titik,
kaya kang diramalake sawetara teori kosmologi.
Yen mangkono, mung ana cacahe detik kang winates
(senajan gedhé banget) saka wektu wiwitan mau nalika
$f(n)$ ditegesi. Saben fungsi kang mung ditegesi
kanggo input winates cacahe bisa diitung: outpute
bisa didaftar ing siji tabel gedhé, utawa dikode
ing siji diagram transisi kahanan mesin Turing
kang gedhé banget.

Kita kerep mligi nggatekake fungsi kang nilaine
menehi wangsulan kanggo pitakon ya/ora. Contone,
pitakon ``apa $n$ wilangan prima?'' digandhengake
karo fungsi
\[
\fn{isprime}(n) = \begin{cases}
  1 & \text{yen $n$ prima}\\
  0 & \text{yen ora.}
  \end{cases}
\]
Kita nyebut pitakon ya/ora \emph{bisa diputusake
kanthi efektif} yen fungsi kang gegayutan lan
nduwé nilai $1/0$ bisa diitung kanthi efektif.

Kanggo mbuktekake sacara matematis yen ana fungsi
kang ora bisa diitung kanthi efektif, utawa masalah
kang ora bisa diputusake kanthi efektif, kita kudu
netepake model komputasi tartamtu banjur nuduhake
yen ana fungsi kang ora bisa diitung utawa masalah
kang ora bisa diputusake dening model iku. Contone,
kita bisa nuduhake yen ora saben fungsi bisa diitung
dening mesin Turing lan ora saben masalah bisa
diputusake dening mesin Turing. Banjur kita bisa
nggunakake tesis Church--Turing kanggo nyimpulake
yen ora mung mesin Turing kang ora cukup kuwat
kanggo ngitung saben fungsi, nanging prosedur
efektif apa wae uga ora bisa.

Kunci kanggo mbuktekake asil negatif kaya mangkono
yaiku kasunyatan yen mesin Turing dhewe bisa diwenehi
nomer. Cara kang paling gampang yaiku ngenumerasi
mesin-mesin iku, upamane kanthi netepake cara tartamtu
kanggo nulis mesin Turing lan programe, banjur
ndhaftar kabeh kanthi runtut. Sawise ngerti yen
iki bisa ditindakake, anane fungsi kang ora bisa
diitung dening mesin Turing dadi akibat saka
petungan kardinalitas kang prasaja: himpunan fungsi
saka $\Nat$ menyang~$\Nat$ (malah saka $\Nat$ menyang
$\{0, 1\}$ wae) iku !!{nonenumerable}, nanging
amarga kabeh mesin Turing bisa dienumerasi, himpunan
fungsi kang bisa diitung dening mesin Turing mung
!!{denumerable}.

Kita uga bisa netepake fungsi lan masalah
\emph{tartamtu} kang bisa dibuktekake ora bisa
diitung lan ora bisa diputusake, miturut urutane.
Salah siji masalah iku diarani \emph{Masalah Mandheg.}
Mesin Turing bisa diterangake kanthi winates
kanthi ndhaptar instruksine. Andharan mesin
Turing kaya mangkono, yaiku program mesin Turing,
mesthi bisa digunakake minangka input kanggo
mesin Turing liyane. Mula kita bisa nliti mesin
Turing kang mutusake pitakon bab mesin Turing liyane.
Salah siji pitakon kang narik kawigaten yaiku:
``Apa mesin Turing kang diwenehake bakal mandheg
ing sawijining wektu yen diwiwiti nganggo input~$n$?''
Apik yen ana mesin Turing kang bisa mutusake pitakon
iki: bayangna minangka mesin pangawas mutu kang
mesthekake yen mesin Turing ora kejebak ing daur
tanpa wates lan sapiturute. Kasunyatan kang narik
kawigaten, kang dibuktekake Turing, yaiku mesin kaya
mangkono ora mungkin ana. Ora ana siji mesin Turing
kang, yen diwiwiti nganggo input kang ngemot andharan
mesin Turing $M$ lan sawijining wilangan~$n$, mesthi
mandheg kanthi output $1$ utawa $0$ miturut apa
$M$ bakal mandheg yen diwiwiti nganggo input $n$
utawa ora.

Yen wis nduwé tuladha masalah tartamtu kang ora bisa
diputusake, kita bisa nggunakake kanggo nuduhake
yen masalah liyane uga ora bisa diputusake. Contone,
salah siji masalah kang misuwur yaiku pitakon,
``Apa !!{formula} orde kapisan~$!A$ sah?''
Ora ana mesin Turing kang, yen diwenehi input
!!{formula} orde kapisan~$!A$, mesthi mandheg
kanthi output $1$ utawa $0$ manut apa $!A$
sah utawa ora. Miturut sejarahe, pitakon golek
prosedur kanggo ngrampungake masalah iki kanthi
efektif mung diarani masalah putusan;
mula kita nyebut masalah putusan iku ora bisa
dirampungake. Turing lan Church mbuktekake asil
iki kanthi mandhiri ing wektu kang meh padha,
mula asil iki uga diarani Teorema Church--Turing.

\end{document}
